\section{Implementacja Bazy Danych oraz Logiki Biznesowej}

\subsection{Implementacja Bazy Danych oraz Logiki Biznesowej}

\subsubsection{Model Aukcji (Auction)}

Model Auction stanowi rdzeń systemu aukcyjnego i zawiera logikę biznesową związaną z zarządzaniem aukcjami.

\begin{longtable}{|p{3.5cm}|p{2.5cm}|p{8cm}|}
  \hline
  \textbf{Pole} & \textbf{Typ}    & \textbf{Opis}                                                                                                                                                 \\ \hline
  \endfirsthead
  \hline
  \textbf{Pole} & \textbf{Typ}    & \textbf{Opis}                                                                                                                                                 \\ \hline
  \endhead

  title         & String          & Tytuł aukcji (3-100 znaków, wymagane)                                                                                                                         \\ \hline
  description   & String          & Opis aukcji (10-5000 znaków, wymagane)                                                                                                                        \\ \hline
  category      & Enum            & 12 kategorii: Sztuka, Antyki, Biżuteria, Monety i Banknoty, Książki, Militaria, Meble, Porcelana i Ceramika, Zegarki, Instrumenty Muzyczne, Elektronika, Inne \\ \hline
  images        & Array[String]   & Tablica URLi obrazów (walidowane URLe)                                                                                                                        \\ \hline
  model3D       & String          & Opcjonalny model 3D (formaty: .obj, .gltf, .glb, .fbx)                                                                                                        \\ \hline
  startingPrice & Number          & Cena wywoławcza (minimalna, wymagana, $\geq 0$)                                                                                                               \\ \hline
  currentPrice  & Number          & Aktualna cena aukcji (wymagana, $\geq 0$)                                                                                                                     \\ \hline
  bidIncrement  & Number          & Minimalne podwyższenie (wymagane, $> 0$, domyślnie: 100)                                                                                                      \\ \hline
  reservePrice  & Number          & Opcjonalna cena minimalna do sprzedaży ($\geq$ startingPrice)                                                                                                 \\ \hline
  buyNowPrice   & Number          & Opcjonalna cena "Kup Teraz" ($>$ startingPrice)                                                                                                               \\ \hline
  seller        & ObjectId        & Referencja do modelu User (wymagane, indeksowane)                                                                                                             \\ \hline
  startTime     & Date            & Data rozpoczęcia aukcji (wymagane, indeksowane)                                                                                                               \\ \hline
  endTime       & Date            & Data zakończenia aukcji (wymagane, $>$ startTime, indeksowane)                                                                                                \\ \hline
  status        & Enum            & Status: draft, active, completed, cancelled (domyślnie: draft, indeksowane)                                                                                   \\ \hline
  winnerId      & ObjectId        & Zwycięzca aukcji (referencja do User)                                                                                                                         \\ \hline
  winningBid    & Number          & Zwycięska oferta                                                                                                                                              \\ \hline
  totalBids     & Number          & Liczba wszystkich licytacji (domyślnie: 0)                                                                                                                    \\ \hline
  watchers      & Array[ObjectId] & Tablica obserwujących aukcję                                                                                                                                  \\ \hline
  watchersCount & Number          & Liczba obserwujących (domyślnie: 0)                                                                                                                           \\ \hline
  views         & Number          & Liczba wyświetleń (domyślnie: 0)                                                                                                                              \\ \hline
  condition     & Enum            & Stan przedmiotu: Nowy, Bardzo dobry, Dobry, Zadowalający, Do renowacji                                                                                        \\ \hline
  featured      & Boolean         & Czy aukcja jest wyróżniona (domyślnie: false)                                                                                                                 \\ \hline
  isDeleted     & Boolean         & Flaga soft delete (domyślnie: false)                                                                                                                          \\ \hline
  timestamps    & Date            & Automatyczne pola: createdAt, updatedAt                                                                                                                       \\ \hline
\end{longtable}

\paragraph{Indeksy Bazy Danych}

System wykorzystuje \textbf{9 strategicznych indeksów} dla optymalizacji zapytań:

\begin{itemize}
  \item \textbf{Złożone:}
        \begin{itemize}
          \item \texttt{seller + status} --- zapytania o aukcje konkretnego sprzedawcy
          \item \texttt{endTime + status} --- wyszukiwanie wygasających aukcji
          \item \texttt{category + status} --- filtrowanie po kategorii
          \item \texttt{featured + status} --- pobieranie wyróżnionych aukcji
        \end{itemize}
  \item \textbf{Pojedyncze:}
        \begin{itemize}
          \item \texttt{currentPrice} --- sortowanie po cenie
          \item \texttt{createdAt} --- sortowanie po dacie utworzenia
        \end{itemize}
  \item \textbf{Tekstowy:}
        \begin{itemize}
          \item \texttt{title + description} --- full-text search w MongoDB
        \end{itemize}
\end{itemize}

\paragraph{Pola Wirtualne}

\begin{itemize}
  \item \texttt{bids} --- relacja wirtualna z modelem Bid (populacja licytacji)
  \item \texttt{timeRemaining} --- oblicza milisekundy pozostałe do końca aukcji
  \item \texttt{isActive} --- zwraca true, jeśli aukcja jest aktywna w danym momencie
  \item \texttt{reserveMet} --- sprawdza, czy osiągnięto cenę rezerwową
\end{itemize}

\paragraph{Metody Instancji}
\begin{enumerate}
  \item incrementViews()
        Inkrementuje licznik wyświetleń aukcji bez uruchamiania walidacji. Wykorzystywane przy każdym otwarciu szczegółów aukcji.
  \item addWatcher(userId)
        Dodaje użytkownika do listy obserwujących (watchlist). Sprawdza duplikaty przed dodaniem. Automatycznie aktualizuje pole \texttt{watchersCount}.
  \item removeWatcher(userId)
        Usuwa użytkownika z listy obserwujących. Wykorzystuje metodę \texttt{filter()} z porównaniem ObjectId przez \texttt{equals()}.
  \item placeBid(bidderId, amount)
        Kluczowa metoda biznesowa do składania licytacji. Przeprowadza  walidację:
        \begin{itemize}
          \item Czy aukcja ma status \texttt{active}
          \item Czy aukcja nie przekroczyła \texttt{endTime}
          \item Czy licytujący nie jest właścicielem aukcji
          \item Czy kwota $\geq$ \texttt{currentPrice + bidIncrement}
        \end{itemize}

        Po walidacji aktualizuje:
        \begin{itemize}
          \item \texttt{currentPrice} = amount
          \item \texttt{winnerId} = bidderId
          \item \texttt{totalBids++}
        \end{itemize}

  \item complete()
        Zamyka aukcję i zmienia status na \texttt{completed}. Jeśli istnieje zwycięzca:
        \begin{itemize}
          \item Ustawia \texttt{winningBid} = \texttt{currentPrice}
          \item Aktualizuje statystyki sprzedawcy: \texttt{stats.totalItemsSold++}
          \item Aktualizuje statystyki zwycięzcy: \texttt{stats.totalAuctionsWon++}
        \end{itemize}
  \item cancel(reason)
        Anuluje aukcję (jeśli status $\neq$ \texttt{completed}). Zmienia status na \texttt{cancelled}.
  \item softDelete()
        Miękkie usunięcie --- ustawia flagę \texttt{isDeleted = true}, zachowując dane w bazie dla celów historycznych i audytu.
\end{enumerate}

\paragraph{Metody Statyczne}
\begin{enumerate}
  \item findActive(options)
        Wyszukuje aktywne aukcje z zaawansowanym filtrowaniem.
        \begin{itemize}
          \item Paginacja: page, limit
          \item Filtrowanie: category, minPrice, maxPrice
          \item Sortowanie: sortBy, sortOrder (domyślnie: endTime, asc
        \end{itemize})
        Automatycznie populuje pole \texttt{seller} (username, rating, avatar).
  \item searchAuctions(query, options)
        Implementuje \textbf{full-text search} wykorzystując indeks tekstowy MongoDB:
        \begin{itemize}
          \item Operator \texttt{\$text: \{ \$search: query \}}
          \item Sortowanie po \texttt{relevance score} (meta: textScore)
          \item Filtrowanie jak w \texttt{findActive()}
          \item Zwraca: aukcje, totalPages, currentPage, total
        \end{itemize}
  \item closeExpiredAuctions()
        Metoda \textbf{batch processing} do automatycznego zamykania aukcji. Wyszukuje wszystkie aukcje gdzie:
        \begin{itemize}
          \item \texttt{status = 'active'}
          \item \texttt{endTime $\leq$ Date.now()}
        \end{itemize}
        Wywołuje metodę \texttt{complete()} na każdej znalezionej aukcji. Zwraca liczbę zamkniętych aukcji. Kompatybilna z cron jobs.
  \item getFeatured(limit)
        Pobiera wyróżnione aukcje (\texttt{featured = true}) o statusie \texttt{active}, sortowane od najnowszych (\texttt{createdAt DESC}).
\end{enumerate}

\newpage

\subsubsection{Model Licytacji (Bid)}

Model Bid reprezentuje pojedynczą licytację w systemie i zawiera logikę walidacji oraz automatycznej aktualizacji powiązanych dokumentów.

\begin{longtable}{|p{3.5cm}|p{2.5cm}|p{8cm}|}
  \hline
  \textbf{Pole} & \textbf{Typ} & \textbf{Opis}                                                    \\ \hline
  \endfirsthead
  \hline
  \textbf{Pole} & \textbf{Typ} & \textbf{Opis}                                                    \\ \hline
  \endhead

  auction       & ObjectId     & Referencja do Auction (wymagane, indeksowane)                    \\ \hline
  bidder        & ObjectId     & Referencja do User (wymagane, indeksowane)                       \\ \hline
  amount        & Number       & Kwota licytacji (wymagane, $\geq 0$)                             \\ \hline
  previousBid   & Number       & Poprzednia cena aukcji (domyślnie: null)                         \\ \hline
  bidType       & Enum         & Typ: manual, auto, buyNow (domyślnie: manual)                    \\ \hline
  isWinning     & Boolean      & Czy to aktualnie wygrywająca licytacja (domyślnie: false)        \\ \hline
  status        & Enum         & Status: active, outbid, won, lost, cancelled (domyślnie: active) \\ \hline
  ipAddress     & String       & Adres IP (do detekcji nadużyć)                                   \\ \hline
  userAgent     & String       & User agent (śledzenie urządzenia)                                \\ \hline
  timestamps    & Date         & createdAt, updatedAt                                             \\ \hline
\end{longtable}

\paragraph{Indeksy Bazy Danych}

System wykorzystuje \textbf{5 indeksów strategicznych}:

\begin{itemize}
  \item \textbf{Złożone:}
        \begin{itemize}
          \item \texttt{auction + bidder} --- sprawdzanie czy użytkownik już licytował
          \item \texttt{auction + amount (desc)} --- szybkie wyszukiwanie najwyższych ofert
          \item \texttt{bidder + createdAt} --- historia licytacji użytkownika
          \item \texttt{auction + isWinning} --- aktualne wygrywające licytacje
        \end{itemize}
  \item \textbf{Pojedyncze:}
        \begin{itemize}
          \item \texttt{createdAt} --- sortowanie chronologiczne
        \end{itemize}
\end{itemize}

\paragraph{Pola Wirtualne}

\begin{itemize}
  \item \texttt{bidIncrease} --- oblicza wzrost kwoty względem poprzedniej licytacji: \texttt{amount - previousBid}
\end{itemize}

\paragraph{Metody Instancji}

\begin{enumerate}
  \item cancel(reason)
        Anuluje aktywną licytację (tylko status \texttt{active}). Ustawia \texttt{status = 'cancelled'} i \texttt{isWinning = false}.

  \item markAsWon()
        Oznacza licytację jako wygraną: \texttt{status = 'won'}, \texttt{isWinning = true}.

  \item markAsLost()
        Oznacza licytację jako przegraną: \texttt{status = 'lost'}, \texttt{isWinning = false}.
\end{enumerate}

\paragraph{Metody Statyczne}
\begin{enumerate}
  \item getAuctionBids(auctionId, options)
        Pobiera historię wszystkich licytacji dla konkretnej aukcji z paginacją (domyślnie: limit 50, sortowanie po \texttt{createdAt DESC}).

  \item getUserBids(userId, options)
        Pobiera historię licytacji konkretnego użytkownika z możliwością filtrowania po statusie i paginacją.

  \item getHighestBid(auctionId)
        Zwraca najwyższą licytację dla danej aukcji (sortowanie po \texttt{amount DESC}, limit 1).

  \item getUserActiveBids(userId)
        Pobiera obecnie wygrywające licytacje użytkownika (\texttt{isWinning = true} AND \texttt{status = 'active'}).

  \item markLostBids(auctionId, winnerId)
        \textbf{Batch update} --- oznacza wszystkie licytacje (oprócz zwycięskiej) jako przegrane po zakończeniu aukcji:
        \begin{lstlisting}
updateMany(
  { auction: auctionId,
    bidder: { $ne: winnerId },
    status: { $in: ['active', 'outbid'] } },
  { status: 'lost', isWinning: false }
)
\end{lstlisting}
\end{enumerate}

\subsection{Model Użytkownika (User)}

Model User zawiera podstawowe informacje o użytkowniku oraz system autentykacji.

\begin{longtable}{|p{3.5cm}|p{2.5cm}|p{8cm}|}
  \hline
  \textbf{Pole}  & \textbf{Typ} & \textbf{Opis}                                                                       \\ \hline
  \endfirsthead
  \hline
  \textbf{Pole}  & \textbf{Typ} & \textbf{Opis}                                                                       \\ \hline
  \endhead

  username       & String       & Nazwa użytkownika (unikalne, indeksowane)                                           \\ \hline
  email          & String       & Adres email (unikalne, indeksowane)                                                 \\ \hline
  password       & String       & Hasło (bcrypt hash, nie zwracane w response)                                        \\ \hline
  firstName      & String       & Imię                                                                                \\ \hline
  lastName       & String       & Nazwisko                                                                            \\ \hline
  phone          & String       & Telefon                                                                             \\ \hline
  avatar         & String       & URL avatara                                                                         \\ \hline
  address        & Object       & Obiekt: street, city, postalCode, country                                           \\ \hline
  rating         & Number       & Średnia ocena użytkownika                                                           \\ \hline
  ratingCount    & Number       & Liczba ocen                                                                         \\ \hline
  accountBalance & Number       & Saldo konta                                                                         \\ \hline
  role           & Enum         & Rola: user, admin                                                                   \\ \hline
  isActive       & Boolean      & Czy konto jest aktywne                                                              \\ \hline
  isVerified     & Boolean      & Czy email zweryfikowany                                                             \\ \hline
  stats          & Object       & Statystyki: totalAuctionsCreated, totalAuctionsWon, totalBidsPlaced, totalItemsSold \\ \hline
  tokens         & Object       & Tokeny: jwt, verification, passwordReset                                            \\ \hline
\end{longtable}

\paragraph{Zabezpieczenia}

\begin{itemize}
  \item Hasło hashowane przez \textbf{bcrypt} (cost factor: 12)
  \item Pole \texttt{password} wyłączone z domyślnych queries (\texttt{select: false})
  \item JWT token z 7-dniowym czasem wygaśnięcia
  \item Weryfikacja aktywności konta (\texttt{isActive = true})
  \item Możliwość blokowania użytkowników przez administratora
\end{itemize}

\subsection{Strategia Walidacji}

System implementuje \textbf{wielopoziomową walidację} danych:

\paragraph{Poziom 1: Schema (Mongoose)}
\begin{itemize}
  \item Typy danych (String, Number, Date, Boolean)
  \item Wymagalność pól (\texttt{required})
  \item Ograniczenia długości (\texttt{minlength}, \texttt{maxlength})
  \item Wartości enum
  \item Custom validators (np. walidacja URL, formatów plików)
\end{itemize}

\paragraph{Poziom 2: Model (Business Logic)}
\begin{itemize}
  \item Metody instancji z logiką biznesową
  \item Middleware hooks (pre/post save)
  \item Sprawdzanie relacji między encjami
  \item Warunki złożone (np. \texttt{buyNowPrice > startingPrice})
\end{itemize}

\paragraph{Poziom 3: Kontroler (Authorization)}
\begin{itemize}
  \item Weryfikacja własności zasobów
  \item Sprawdzanie ról użytkownika (user/admin)
  \item Kontrola dostępu do operacji
\end{itemize}

\subsection{Integralność Danych}

\paragraph{Relacje i Referencje}
\begin{itemize}
  \item \textbf{User $\rightarrow$ Auction} (1:N) --- pole \texttt{seller}
  \item \textbf{User $\rightarrow$ Bid} (1:N) --- pole \texttt{bidder}
  \item \textbf{Auction $\rightarrow$ Bid} (1:N) --- pole \texttt{auction}
  \item \textbf{Auction $\leftrightarrow$ User} (N:M) --- tablica \texttt{watchers}
\end{itemize}

\newpage

\subsection{Implementacja Ścieżek API}

\subsubsection{Architektura REST API}

API implementuje standardy \textbf{RESTful} z następującymi konwencjami:
\begin{itemize}
  \item \textbf{GET} --- pobieranie zasobów
  \item \textbf{POST} --- tworzenie zasobów
  \item \textbf{PATCH} --- częściowa aktualizacja zasobów
  \item \textbf{DELETE} --- usuwanie zasobów
\end{itemize}

\subsubsection{Middleware Autentykacji}

System wykorzystuje dwa główne middleware do kontroli dostępu.

\paragraph{Middleware protect}

\textbf{Funkcjonalność:}
\begin{enumerate}
  \item Ekstrakcja JWT z nagłówka \texttt{Authorization: Bearer <token|}
  \item Weryfikacja podpisu tokena (klucz: \texttt{JWT\_SECRET})
  \item Pobranie użytkownika z bazy danych po \texttt{userId} z tokena
  \item Sprawdzenie czy konto aktywne (\texttt{isActive = true})
  \item Walidacja czy hasło nie zostało zmienione po wydaniu tokena
  \item Dołączenie obiektu \texttt{user} do \texttt{req.user}
  \item Zwrot \texttt{401 Unauthorized} dla brakujących/nieprawidłowych tokenów
\end{enumerate}

\paragraph{Middleware restrictTo(...roles)}

\textbf{Funkcjonalność:}
\begin{itemize}
  \item Kontrola dostępu oparta na rolach (\textbf{RBAC --- Role-Based Access Control})
  \item Akceptuje zmienną liczbę dozwolonych ról jako parametry
  \item Zwrot \texttt{403 Forbidden} jeśli rola użytkownika nie znajduje się na liście
  \item Używane dla endpointów wymagających uprawnień administratora
\end{itemize}

\textbf{Przykład użycia:}
\begin{lstlisting}
router.post('/admin/close-expired',
  protect,
  restrictTo('admin'),
  AuctionController.closeExpiredAuctions
);
\end{lstlisting}

\subsubsection{API Aukcji}

\paragraph{Endpointy Publiczne (bez autoryzacji)}

\begin{enumerate}
  \item \textbf{GET /api/auctions}

        \textbf{Opis:} Pobieranie listy aukcji z filtrowaniem i paginacją.

        \textbf{Query parametry:}
        \begin{itemize}
          \item \texttt{page} --- numer strony (domyślnie: 1)
          \item \texttt{limit} --- liczba wyników na stronę (domyślnie: 10)
          \item \texttt{category} --- filtrowanie po kategorii
          \item \texttt{minPrice, maxPrice} --- zakres cenowy
          \item \texttt{sortBy} --- pole sortowania (createdAt, currentPrice, endTime)
          \item \texttt{sortOrder} --- kierunek (asc/desc)
          \item \texttt{status} --- status aukcji (draft, active, completed, cancelled)
        \end{itemize}

        \textbf{Response:}
        \begin{lstlisting}
{
  success: true,
  results: Number,
  totalPages: Number,
  currentPage: Number,
  auctions: [Auction]
}
\end{lstlisting}

  \item \textbf{GET /api/auctions/search}

        \textbf{Opis:} Wyszukiwanie full-text w tytułach i opisach aukcji.

        \textbf{Query parametry:}
        \begin{itemize}
          \item \texttt{q} --- zapytanie tekstowe (wymagane)
          \item \texttt{page, limit} --- paginacja
          \item \texttt{category, minPrice, maxPrice} --- filtrowanie
          \item \texttt{status} --- status aukcji
          \item \texttt{sortBy} --- domyślnie: relevance score
        \end{itemize}

        \textbf{Implementacja:} Wykorzystuje indeks tekstowy MongoDB (\texttt{\$text: \{ \$search: query \}}) z sortowaniem po \texttt{textScore}.

  \item \textbf{GET /api/auctions/:id}

        \textbf{Opis:} Szczegóły pojedynczej aukcji. Automatycznie inkrementuje licznik \texttt{views}.

        \textbf{Response:} Aukcja z populowanymi polami \texttt{seller} i \texttt{winner}.

        \textbf{Błędy:}
        \begin{itemize}
          \item \texttt{404} --- Aukcja nie znaleziona lub usunięta
        \end{itemize}

  \item \textbf{GET /api/auctions/:id/bids}

        \textbf{Opis:} Historia licytacji dla konkretnej aukcji.

        \textbf{Query parametry:} page, limit, sortBy, sortOrder
\end{enumerate}

\paragraph{Endpointy Chronione (wymagana autentykacja)}

\begin{enumerate}
  \item \textbf{POST /api/auctions}

        \textbf{Middleware:} \texttt{protect}

        \textbf{Opis:} Utworzenie nowej aukcji.

        \textbf{Body (wymagane pola):}
        \begin{lstlisting}
{
  title: String (3-100 znakow),
  description: String (10-5000 znakow),
  category: String (jedna z 12 kategorii),
  startingPrice: Number (> 0),
  bidIncrement: Number (> 0),
  startTime: Date (>= now),
  endTime: Date (> startTime),
  condition: String (enum)
}
\end{lstlisting}

        \textbf{Pola opcjonalne:}
        \begin{itemize}
          \item \texttt{images} --- tablica URLi obrazów
          \item \texttt{model3D} --- URL modelu 3D
          \item \texttt{reservePrice} --- cena minimalna
          \item \texttt{buyNowPrice} --- cena "Kup Teraz"
        \end{itemize}

        \textbf{Automatyzacja:}
        \begin{itemize}
          \item \texttt{seller = req.user.\_id}
          \item \texttt{user.stats.totalAuctionsCreated++}
        \end{itemize}

        \textbf{Błędy:}
        \begin{itemize}
          \item \texttt{400} --- Błędy walidacji
          \item \texttt{401} --- Brak autoryzacji
        \end{itemize}

  \item \textbf{PATCH /api/auctions/:id}

        \textbf{Middleware:} \texttt{protect}

        \textbf{Opis:} Aktualizacja aukcji (tylko właściciel lub admin).

        \textbf{Ograniczenia edycji:}
        \begin{itemize}
          \item \textbf{Aukcje draft:} można edytować wszystkie pola
          \item \textbf{Aukcje active bez bidów:} można edytować description, images
          \item \textbf{Aukcje active z bidami:} brak możliwości edycji
          \item \textbf{Aukcje completed/cancelled:} brak możliwości edycji
        \end{itemize}

        \textbf{Błędy:}
        \begin{itemize}
          \item \texttt{403} --- Brak uprawnień (nie właściciel)
          \item \texttt{400} --- Edycja aktywnej aukcji z bidami zabroniona
          \item \texttt{404} --- Aukcja nie znaleziona
        \end{itemize}

  \item \textbf{DELETE /api/auctions/:id}

        \textbf{Middleware:} \texttt{protect}

        \textbf{Opis:} Soft delete aukcji (tylko właściciel lub admin).

        \textbf{Ograniczenia:}
        \begin{itemize}
          \item Nie można usunąć active auction z bidami
        \end{itemize}

  \item \textbf{POST /api/auctions/:id/bid}

        \textbf{Middleware:} \texttt{protect}

        \textbf{Opis:} Złożenie licytacji.

        \textbf{Body:}
        \begin{lstlisting}
{
  amount: Number (>= currentPrice + bidIncrement)
}
\end{lstlisting}

        \textbf{Walidacja:}
        \begin{itemize}
          \item Aukcja musi być aktywna (\texttt{status = active}, \texttt{startTime $\leq$ now $\leq$ endTime})
          \item Kwota $\geq$ \texttt{currentPrice + bidIncrement}
          \item Nie można licytować własnej aukcji
          \item Użytkownik nie może przelicytować samego siebie
        \end{itemize}

        \textbf{Automatyzacja (post-save hook):}
        \begin{itemize}
          \item \texttt{auction.currentPrice = amount}
          \item \texttt{auction.winnerId = bidderId}
          \item \texttt{auction.totalBids++}
          \item Poprzednie bidy: \texttt{status = outbid}
          \item \texttt{user.stats.totalBidsPlaced++}
        \end{itemize}

        \textbf{Błędy:}
        \begin{itemize}
          \item \texttt{400} --- Kwota za niska
          \item \texttt{403} --- Licytacja własnej aukcji
          \item \texttt{404} --- Aukcja nie znaleziona
          \item \texttt{400} --- Aukcja nieaktywna
        \end{itemize}

  \item \textbf{GET /api/auctions/me/auctions}

        \textbf{Middleware:} \texttt{protect}

        \textbf{Opis:} Lista aukcji utworzonych przez użytkownika z możliwością filtrowania po statusie.

  \item \textbf{POST /api/auctions/:id/watch}

        \textbf{Middleware:} \texttt{protect}

        \textbf{Opis:} Dodanie aukcji do watchlist.

        \textbf{Automatyzacja:}
        \begin{itemize}
          \item \texttt{auction.watchers.push(userId)}
          \item \texttt{auction.watchersCount++}
        \end{itemize}

  \item \textbf{DELETE /api/auctions/:id/watch}

        \textbf{Middleware:} \texttt{protect}

        \textbf{Opis:} Usunięcie aukcji z watchlist.
\end{enumerate}

\subsubsection{API Licytacji}

\paragraph{Endpointy Publiczne}

\begin{enumerate}
  \item \textbf{GET /api/bids/auction/:auctionId/stats}

        \textbf{Opis:} Statystyki agregowane dla aukcji (MongoDB aggregation pipeline).

        \textbf{Response:}
        \begin{lstlisting}
{
  success: true,
  stats: {
    totalBids: Number,
    uniqueBidders: Number,
    averageBid: Number,
    minBid: Number,
    maxBid: Number
  }
}
\end{lstlisting}
\end{enumerate}

\paragraph{Endpointy Chronione}

\begin{enumerate}
  \item \textbf{GET /api/bids/me}

        \textbf{Middleware:} \texttt{protect}

        \textbf{Opis:} Historia licytacji użytkownika z możliwością filtrowania po statusie.

  \item \textbf{PATCH /api/bids/:id/cancel}

        \textbf{Middleware:} \texttt{protect, restrictTo('admin')}

        \textbf{Opis:} Anulowanie licytacji (tylko admin).

        \textbf{Body:}
        \begin{lstlisting}
{
  reason: String (opcjonalne)
}
\end{lstlisting}

        \textbf{Ograniczenia:}
        \begin{itemize}
          \item Tylko admin może anulować
          \item Nie można anulować bidów ze statusem won/lost
        \end{itemize}

        \textbf{Błędy:}
        \begin{itemize}
          \item \texttt{403} --- Brak uprawnień administratora
          \item \texttt{400} --- Bid nie może być anulowany (completed)
        \end{itemize}
\end{enumerate}

\subsubsection{Kody Statusów HTTP}

\begin{longtable}{|p{2cm}|p{4cm}|p{8cm}|}
  \hline
  \textbf{Kod} & \textbf{Nazwa}        & \textbf{Użycie}                             \\ \hline
  \endfirsthead
  \hline
  \textbf{Kod} & \textbf{Nazwa}        & \textbf{Użycie}                             \\ \hline
  \endhead

  200          & OK                    & Sukces dla GET/PATCH                        \\ \hline
  201          & Created               & Sukces dla POST (utworzenie zasobu)         \\ \hline
  400          & Bad Request           & Błędy walidacji lub logiki biznesowej       \\ \hline
  401          & Unauthorized          & Brak lub nieprawidłowy token JWT            \\ \hline
  403          & Forbidden             & Brak uprawnień do operacji (rola, własność) \\ \hline
  404          & Not Found             & Zasób nie istnieje                          \\ \hline
  500          & Internal Server Error & Błędy serwera                               \\ \hline
\end{longtable}

\subsubsection{Format Odpowiedzi}

\paragraph{Sukces}
\begin{lstlisting}
{
  "success": true,
  "data": { ... },
  "results": 10,
  "totalPages": 5,
  "currentPage": 1
}
\end{lstlisting}

\paragraph{Błąd}
\begin{lstlisting}
{
  "success": false,
  "error": {
    "message": "Opis bledu",
    "statusCode": 400
  }
}
\end{lstlisting}

W trybie deweloperskim błędy zawierają dodatkowo \texttt{stack trace}.

\subsubsection{Przepływ Danych --- Kluczowe Scenariusze}

\paragraph{Scenariusz 1: Złożenie Licytacji}

\begin{enumerate}
  \item Użytkownik: \texttt{POST /api/auctions/:id/bid \{ amount \}}
  \item \texttt{AuctionController.placeBid()} --- walidacja żądania HTTP
  \item \texttt{Auction.placeBid(bidderId, amount)} --- walidacja biznesowa:
        \begin{itemize}
          \item Czy aukcja aktywna?
          \item Czy kwota wystarczająca?
          \item Czy nie własna aukcja?
        \end{itemize}
  \item Utworzenie dokumentu Bid
  \item \textbf{Bid pre-save hook} --- ponowna walidacja na poziomie modelu
  \item \textbf{Bid post-save hook} --- automatyzacja:
        \begin{itemize}
          \item \texttt{Auction.currentPrice = amount}
          \item \texttt{Auction.winnerId = bidderId}
          \item Poprzednie bidy: \texttt{status = outbid}, \texttt{isWinning = false}
          \item \texttt{User.stats.totalBidsPlaced++}
        \end{itemize}
  \item Zwrot zaaktualizowanej licytacji do klienta
\end{enumerate}

\paragraph{Scenariusz 2: Zakończenie Aukcji}

\begin{enumerate}
  \item Wywołanie \texttt{Auction.complete()} (ręcznie lub przez cron)
  \item Zmiana \texttt{status = 'completed'}
  \item Jeśli istnieje \texttt{winnerId}:
        \begin{itemize}
          \item \texttt{winningBid = currentPrice}
          \item \texttt{Seller.stats.totalItemsSold++}
          \item \texttt{Winner.stats.totalAuctionsWon++}
        \end{itemize}
  \item Zapis do bazy danych
  \item Wywołanie \texttt{Bid.markLostBids()} dla pozostałych licytacji
\end{enumerate}

\paragraph{Scenariusz 3: Wyszukiwanie Full-Text}

\begin{enumerate}
  \item Użytkownik: \texttt{GET /api/auctions/search?q=query}
  \item \texttt{AuctionController.searchAuctions()} --- ekstrakcja parametrów
  \item \texttt{Auction.searchAuctions(query)} --- użycie indeksu tekstowego MongoDB
  \item Sortowanie wyników po relevance score (\texttt{\$meta: 'textScore'})
  \item Paginacja i zwrot do klienta
\end{enumerate}
